% ===========================================================================
%  E12 — Confiança: privacidade, segurança e operação
% ===========================================================================
\chapter{E12 \textperiodcentered{} Confiança: privacidade, segurança e operação}
\label{ep:e12}

\begin{fichaepico}
\begin{campos}
  \item[Objetivo] A rede, a escola e a família podem confiar na plataforma:
    ela protege o dado de criança, aguenta a manhã de aula, e a equipe sabe
    quando algo quebra.
  \item[Para quem] Secretaria, encarregado de dados da rede, família --- e,
    indiretamente, todos os papéis.
  \item[Valor] São as histórias que nenhuma tela mostra, mas que um edital, um
    jurídico ou um incidente cobram. Várias delas bloqueiam a entrada em
    produção.
  \item[Estado] \parcial{} --- isolamento e proteção de menores prontos;
    teste de carga, monitoramento em produção, retenção e acessibilidade
    pendentes.
\end{campos}
\end{fichaepico}

\begin{tabularx}{\linewidth}{@{}L{16mm}X L{20mm}@{}}
\toprule
\headrow \thd{ID} & \thd{Feature} & \thd{Estado}\\
\midrule
\id{F12.1} & Uma escola nunca vê outra & \pronto\\
\zebra \id{F12.2} & Dado de criança protegido & \pronto\\
\id{F12.3} & Guarda e eliminação de dados & \afazer\\
\zebra \id{F12.4} & Aguentar as 8h da manhã & \parcial\\
\id{F12.5} & A equipe sabe quando algo quebra & \parcial\\
\zebra \id{F12.6} & Acessibilidade & \afazer\\
\id{F12.7} & Português do Brasil & \pronto\\
\bottomrule
\end{tabularx}

% ---------------------------------------------------------------------------
\feature{F12.1}{Uma escola nunca vê outra}{\pronto}

\begin{historia}{H12.1.1}{Isolamento garantido entre escolas}{\pronto}
\hu{secretaria}{garantia de que uma escola nunca vê dado de criança de outra
escola}{atender a LGPD e o edital}
\aceite
\begin{itemize}
  \item A separação é garantida na própria base de dados, não só nas telas:
    vale mesmo que haja um erro de programação.
  \item As únicas exceções são as configurações da rede (disponibilidade,
    limites de IA) e os cadastros do site, que não são dado de estudante.
  \item Existe teste automatizado que tenta atravessar escolas e falha.
  \item Cada rede tem instalação própria: um incidente numa rede não atinge
    outra.
\end{itemize}
\end{historia}

% ---------------------------------------------------------------------------
\feature{F12.2}{Dado de criança protegido}{\pronto}

\begin{historia}{H12.2.1}{Nada que identifique a criança onde não precisa}{\pronto}
\hu{encarregado de dados da rede}{que os registros de uso e os relatórios não
identifiquem a criança sem necessidade}{reduzir o risco de vazamento}
\aceite
\begin{itemize}
  \item Os registros de uso carregam um identificador opaco, nunca nome ou
    e-mail; o CPF nunca aparece em registro técnico.
  \item Os relatórios de gestão são agregados; a linha de uma criança é só do
    professor dela.
  \item Onde a missão pede texto, há limite de tamanho, e dado pessoal é
    recusado (Manifesto da IA).
  \item Na Vila do Lume o nome da criança nunca é gravado.
  \item A exceção declarada é A Seleção, em que um colega lê o que a criança
    publicou: no Mural, o colega com quem a publicação foi pareada vê os três
    posts, a legenda e o primeiro nome dela (ou o nome da equipe). Por isso a
    legenda recusa link, telefone, e-mail, marcação de pessoa e palavrão, e a
    Política de Privacidade e os Termos de Uso dizem isso.
\end{itemize}
\pendencia{uma ferramenta de gravação de sessão de terceiros é usada nas telas
e precisa constar no termo de tratamento de dados de cada rede --- ou ser
retirada das telas de menores.}
\end{historia}

% ---------------------------------------------------------------------------
\feature{F12.3}{Guarda e eliminação de dados}{\afazer}

\begin{historia}{H12.3.1}{Prazos de guarda definidos e cumpridos}{\afazer}
\hu{encarregado de dados da rede}{saber por quanto tempo cada dado de criança
é guardado, e que ele seja apagado no prazo}{cumprir a LGPD e responder à
família}
\aceiteproposta
\begin{itemize}
  \item Um prazo definido para cada tipo de dado (resultado de missão, texto
    livre, registros de uso).
  \item Uma rotina automática apaga o que passou do prazo.
  \item A Política de Privacidade passa a declarar esses prazos.
\end{itemize}
\pendencia{decisão de produto e jurídica anterior ao desenvolvimento (\id{D09}).}
\end{historia}

% ---------------------------------------------------------------------------
\feature{F12.4}{Aguentar as 8h da manhã}{\parcial}

\begin{historia}{H12.4.1}{A rede inteira entrando ao mesmo tempo}{\parcial}
\hu{secretaria}{que a plataforma aguente a rede inteira começando a aula ao
mesmo tempo}{que nenhuma turma perca a aula às 8h}
\aceite
\begin{itemize}
  \item 2 a 4 mil pessoas entrando e abrindo missão em 10 minutos, sem erro e
    sem lentidão perceptível.
  \item Inclui turmas inteiras conversando com a IA ao mesmo tempo (unooGPT).
\end{itemize}
\pendencia{a plataforma está configurada para isso, mas o teste de carga
nunca foi feito. É o maior risco aberto antes de novembro.}
\end{historia}

% ---------------------------------------------------------------------------
\feature{F12.5}{A equipe sabe quando algo quebra}{\parcial}

\begin{historia}{H12.5.1}{Saber do problema antes da escola ligar}{\parcial}
\hu{equipe da Unooria}{ser avisada quando a plataforma falha}{resolver antes que
a escola perceba}
\aceite
\begin{itemize}
  \item Meta: 99,5\% das ações sem erro.
  \item Verificações automáticas a cada minuto, de três regiões; alertas por
    e-mail; registro dos erros que acontecem no navegador de quem usa,
    inclusive na tela de entrada.
  \item Nenhum dado pessoal nesses registros.
\end{itemize}
\pendencia{tudo isso está funcionando só no ambiente de desenvolvimento;
falta levar para homologação e produção, testar o canal de e-mail e ter um
alerta de custo.}
\end{historia}

% ---------------------------------------------------------------------------
\feature{F12.6}{Acessibilidade}{\afazer}

\begin{historia}{H12.6.1}{Usar a plataforma com deficiência}{\afazer}
\hu{estudante com deficiência visual ou motora}{usar a plataforma com leitor
de tela e teclado}{participar da aula como os colegas}
\aceiteproposta
\begin{itemize}
  \item Contraste, foco visível, navegação por teclado e leitor de tela
    verificados em todas as telas da plataforma e nas missões.
  \item A transcrição de todo áudio continua sempre visível (isso já existe).
\end{itemize}
\pendencia{hoje não há verificação sistemática; a correção deve ser feita no
componente visual compartilhado, para valer em todas as telas de uma vez.
Até lá, não é possível prometer acessibilidade a uma rede.}
\end{historia}

% ---------------------------------------------------------------------------
\feature{F12.7}{Português do Brasil}{\pronto}

\begin{historia}{H12.7.1}{Tudo na minha língua}{\pronto}
\hu{qualquer pessoa}{que toda a plataforma fale português do
Brasil}{não tropeçar em termo estrangeiro}
\aceite
\begin{itemize}
  \item Toda a interface da plataforma está em português do Brasil.
  \item Dentro do aplicativo imitado, ficam os termos do aplicativo real
    (Feed, Stories, Salvos), e os nomes Unoovista, Instanooria e unooGPT não se traduzem.
  \item \enquote{fake news}, \enquote{cyberbullying} e \enquote{print} ficam
    como estão, por decisão.
  \item Datas no formato brasileiro e no fuso da rede.
\end{itemize}
\end{historia}
